\subsection{Weyr-Fitting Decomposition}

Let A be a ring, M an A-module, and u an endomorphism of M. For any integer \(p \geqslant 0\) , we denote the kernel of \(u^{p}\) by \(N_{p}\) . The sequence of submodules \((\mathrm{N}_{p})\) is increasing, and its union is a submodule \(N_{\infty}\) of M that is stable under u. For every integer \(p \geqslant 0\) , we have \(\mathrm{N}_{p+1} = \overline{u}^{-1}(\mathrm{N}_{p})\) , and the relation \(N_{p} = N_{p+1}\) therefore implies \(N_{p+1} = N_{p+2}\) . Consequently, either the sequence \((\mathrm{N}_{p})\) is strictly increasing, or there exists an integer \(p \geqslant 0\) such that \(N_{0}, \ldots, N_{p}\) are distinct and \(N_{p} = N_{\infty}\) .

For any integer \(q \geqslant 0\) , denote the image of \(u^{q}\) by \(I_{q}\) . The sequence of submodules ( \(I_{q}\) ) is decreasing, and its intersection is a submodule \(I_{\infty}\) of M that is stable under u. For every integer \(q \geqslant 0\) , we have \(u(\mathrm{I}_{q}) = \mathrm{I}_{q+1}\) , and the relation \(I_{q} = I_{q+1}\) therefore implies \(I_{q+1} = I_{q+2}\) . Consequently, either the sequence ( \(I_{q}\) ) is strictly decreasing, or there exists an integer \(q \geqslant 0\) such that \(I_{0}, \ldots, I_{q}\) are distinct and \(I_{q} = I_{\infty}\) .

PROPOSITION 2. --- a) Suppose that the sequence \((\mathrm{N}_{p})\) is stationary. Then we have \(N_{\infty} \cap I_{\infty} = 0\) , the restriction of u to \(I_{\infty}\) is injective, and u induces a nilpotent endomorphism of \(N_{\infty}\) .

\begin{enumerate}
\def\labelenumi{\alph{enumi})}
\setcounter{enumi}{1}
\item
  Suppose that the sequence \((\mathrm{I}_p)\) is stationary. Then we have \(\mathbf{M} = \mathrm{N}_{\infty} + \mathrm{I}_{\infty}\) and \(u(\mathrm{I}_{\infty}) = \mathrm{I}_{\infty}\) .
\item
  (``Weyr--Fitting decomposition''---sometimes called the ``Fitting decomposition'') Suppose that the sequences \((\mathrm{N}_{p})\) and \((\mathrm{I}_{p})\) are stationary. Then M is a direct sum of submodules \(N_{\infty}\) and \(I_{\infty}\) that are stable under u, and u induces a nilpotent endomorphism of \(N_{\infty}\) and an automorphism of \(I_{\infty}\) .
\end{enumerate}

Let p be a natural number such that \(N_{p} = N_{\infty}\) , and let \(v = u^{p}\) . By construction, v and \(v^{2}\) have the same kernel \(N_{\infty}\) , and \(I_{p}\) is the image of v. For x in \(N_{\infty} \cap I_{p}\) , there exists a \(y \in M\) such that \(x = v(y)\) and therefore \(v^{2}(y) = v(x) = 0\) ; hence, we have \(y \in N_{\infty}\) and consequently x = 0. In particular, we have \(N_{\infty} \cap I_{\infty} = 0\) . Since the kernel \(N_{1}\) of u is contained in \(N_{\infty}\) , the restriction of u to \(I_{\infty}\) is injective; on the other hand, we have \(u^{p}(N_{\infty}) = 0\) . This proves a).

Let \(q\) be a natural number such that \(\mathrm{I}_q = \mathrm{I}_{\infty}\) , and let \(w = u^q\) . Then \(w\) and \(w^2\) have the same image \(\mathrm{I}_{\infty}\) , and \(\mathrm{N}_q\) is the kernel of \(w\) . Let \(x \in \mathrm{M}\) . We have \(w(x) \in \mathrm{I}_{\infty}\) , so there exists a \(y \in \mathrm{M}\) such that \(w(x) = w^2(y)\) ; we then have \(x - w(y) \in \mathrm{N}_q\) , whence \(\mathrm{M} = \mathrm{N}_q + \mathrm{I}_{\infty}\) and a fortiori \(\mathrm{M} = \mathrm{N}_{\infty} + \mathrm{I}_{\infty}\) . We have \(u(\mathrm{I}_{\infty}) = u(\mathrm{I}_q) = \mathrm{I}_{q+1} = \mathrm{I}_{\infty}\) . This proves b).

Assertion c) follows immediately from a) and b).

Remarks. --- 1) Let p be an integer such that \(N_{p} = N_{p+1}\) ; the proof given above shows that \(N_{\infty} \cap I_{p} = 0\) , and the restriction of u to \(I_{p}\) is injective. Likewise, let q be an integer such that \(I_{q} = I_{q+1}\) ; then we have \(N_{q} + I_{\infty} = M\) , and the endomorphism of \(M/N_{q}\) deduced from u by passing to the quotient is surjective.

\begin{enumerate}
\def\labelenumi{\arabic{enumi})}
\setcounter{enumi}{1}
\tightlist
\item
  Suppose that M is a direct sum of two submodules N and I that are stable under u and that u induces a nilpotent endomorphism \(u_{N}\) of N and an automorphism of I. We then have \(N_{\infty} = N\) and \(I_{\infty} = I\) , and the sequences \((N_{p})\) and \((I_{p})\) are stationary. Moreover, the following integers are equal:
\end{enumerate}

\(\alpha)\) the least integer \(p \geqslant 0\) such that \(\mathrm{N}_p = \mathrm{N}_{\infty}\) ,

\(\beta)\) the least integer \(q \geqslant 0\) such that \(\mathrm{I}_q = \mathrm{I}_{\infty}\) ,

\(\gamma)\) the least integer \(r \geqslant 0\) such that \((u_{\mathrm{N}})^{r} = 0\) .

\begin{enumerate}
\def\labelenumi{\arabic{enumi})}
\setcounter{enumi}{2}
\tightlist
\item
  The assumption of assertion a) is satisfied if the A-module M is Noetherian; the assumption of assertion b) is satisfied if M is Artinian; by Proposition 1 of VIII, p.~2, the assumption of assertion c) is satisfied if M has finite length.
\end{enumerate}

\textbf{COROLLARY 1. --- Let A be a ring, and let M be an A-module.}

\begin{enumerate}
\def\labelenumi{\alph{enumi})}
\item
  If the module M is Noetherian, then every surjective endomorphism of M is bijective.
\item
  If the module M is Artinian, then every injective endomorphism of M is bijective.
\item
  If the module M has finite length, then every injective or surjective endomorphism of M is bijective.
\item
  If the ring A is commutative and the A-module M is finitely generated, then every surjective endomorphism of M is bijective.
\end{enumerate}

Let u be an endomorphism of the A-module M. We use the notation introduced at the beginning of this subsection. If the endomorphism u is surjective, then we have \(I_{\infty} = M\) . Assertion a) then follows from Proposition 2, a) and Remark 3. Likewise, if the endomorphism u is injective, we have \(N_{\infty} = 0\) . Assertion b) therefore follows from Proposition 2, b) and Remark 3. Assertion c) follows immediately from a) and b).

Now, suppose that the ring A is commutative, the A-module M is finitely generated, and the endomorphism u is surjective. Let us prove that u is injective. Let x be an element of M such that \(u(x) = 0\) . Choose a finite generating family \((x_{i})_{i \in \mathrm{I}}\) of the A-module M and, for every \(i \in I\) , an element \(y_{i}\) of M such that \(u(y_{i}) = x_{i}\) . There exist families \((a_{i})_{i \in \mathrm{I}}, (b_{ij})_{(i,j) \in \mathrm{I} \times \mathrm{I}}\) , and \((c_{ij})_{(i,j) \in \mathrm{I} \times \mathrm{I}}\) of elements of A such that we have

\[
x = \sum_ {i \in \mathrm{I}} a _ {i} x _ {i}, \quad y _ {j} = \sum_ {i \in \mathrm{I}} b _ {i j} x _ {i}, \quad u (x _ {j}) = \sum_ {i \in \mathrm{I}} c _ {i j} x _ {i}
\]

for every \(j \in I\) . Let \(A'\) be a Noetherian subring of A containing the elements \(a_i\) , \(b_{ij}\) , and \(c_{ij}\) (VIII, p.~12, Corollary 3). Let \(M'\) be the \(A'\) -submodule of M generated by the family \((x_i)_{i \in I}\) . We have \(u(x_j) \in M'\) , \(y_j \in M'\) , and \(u(y_j) = x_j\) for every \(j \in I\) ; hence u defines, by restriction, a surjective endomorphism \(u'\) of the \(A'\) -module \(M'\) . Since the ring \(A'\) is Noetherian, the finitely generated \(A'\) -module \(M'\) is Noetherian (VIII, p.~7, Proposition 4 a)). By a), the endomorphism \(u'\) of \(M'\) is bijective. By construction, x belongs to \(M'\) , and we have \(u'(x) = u(x) = 0\) . We therefore have x = 0, which proves d).

COROLLARY 2. --- In a left Noetherian ring, every left or right invertible element is invertible.

Indeed, consider elements \(x\) , \(y\) of a left Noetherian ring \(A\) such that \(xy = 1\) . Denote by \(\pmb{\delta}(x)\) and \(\pmb{\delta}(y)\) , respectively, the endomorphisms \(a \mapsto ax\) and \(a \mapsto ay\) of the A-module \(A_s\) . We have \(\pmb{\delta}(y) \circ \pmb{\delta}(x) = 1_{A_s}\) ; hence \(\pmb{\delta}(y)\) is surjective. By Corollary 1, a), \(\pmb{\delta}(y)\) is bijective. The endomorphism \(\pmb{\delta}(x)\) is then the inverse bijection of \(\pmb{\delta}(y)\) , and we have \(yx = (\pmb{\delta}(x) \circ \pmb{\delta}(y))(1) = 1\) . The corollary follows.
% semantic-fragments: legacy-input-seam
\relax{}
